\chapter{Tổng kết}

\section{So sánh các nhóm bài toán}

Các bài toán đã hoàn thiện trong báo cáo có thể được chia thành ba nhóm lớn. Nhóm thứ nhất tập trung vào tài nguyên hữu hạn như phòng đọc sách và bãi đỗ xe. Nhóm thứ hai tập trung vào ghép nhóm, phối hợp điều kiện và đánh thức đúng tiến trình như River Crossing, Cigarette Smokers, Santa Claus và Child Care. Nhóm thứ ba mô phỏng dịch vụ thực tế có nhiều pha hoạt động như Elevator, Bear and Honeybees và Senate Bus.

Điểm chung của các bài toán là đều yêu cầu xác định đúng tài nguyên găng, trạng thái dùng chung và điều kiện đánh thức. Một lời giải tốt không chỉ ngăn race condition mà còn phải tránh deadlock, hạn chế starvation và đảm bảo hệ thống vẫn có progress trong các tình huống cạnh tranh tài nguyên.

\section{Kết luận}

Qua các nhóm bài toán đã phân tích, có thể thấy mutex, semaphore, monitor và biến điều kiện không chỉ là công cụ lập trình mà còn là cách mô hình hóa quan hệ phối hợp giữa các tiến trình. Khi thiết kế lời giải đồng bộ, cần bắt đầu từ bất biến an toàn của hệ thống, sau đó xác định chính xác tiến trình nào được chờ, tiến trình nào có quyền đánh thức và tài nguyên nào phải được bảo vệ trong vùng găng.
